\begin{table}[htbp]
\centering
\caption{Hit identification by edgeR. Shown compounds are for FDR\textless{}0.05 and sorted by logFC to identify the top hits.}
\label{tab:table12}
\begin{tabular}{lllllll}
\toprule
code\_1 & code\_2 & logFC & logCPM & LR & PValue & FDR \\
\midrule
94 & 472 & 9.76658103 & 4.65399379 & 140.973181 & 1.63E-32 & 2.99E-29 \\
94 & 337 & 9.52452422 & 4.41931582 & 128.115071 & 1.06E-29 & 9.38E-27 \\
250 & 250 & 9.50138143 & 4.39920396 & 136.048263 & 1.95E-31 & 2.62E-28 \\
250 & 523 & 9.46551754 & 4.36448577 & 134.035088 & 5.37E-31 & 6.34E-28 \\
94 & 364 & 9.29500646 & 4.19806216 & 117.172728 & 2.63E-27 & 1.35E-24 \\
250 & 208 & 9.2845154 & 5.99659857 & 173.194896 & 1.48E-39 & 2.05E-35 \\
250 & 337 & 9.06687397 & 3.98375428 & 115.539372 & 6.00E-27 & 2.80E-24 \\
94 & 341 & 8.99590166 & 3.91723403 & 113.847311 & 1.41E-26 & 6.00E-24 \\
250 & 414 & 8.99240694 & 3.90981012 & 105.694583 & 8.60E-25 & 2.69E-22 \\
250 & 96 & 8.92430743 & 8.02773969 & 203.793994 & 3.10E-46 & 2.36E-41Phase 2: Chemical library enumeration and analysis • TIMING 10 min - 2 h This phase generates all possible chemical structures from your building block combinations and calculates molecular properties for downstream analysis. Library enumeration requires advanced knowledge of the SMARTS/SMIRKS language. While we provide common reaction patterns observed in DNA encoded libraries, library specific tweaks are often necessary. Step 4 \textbar{} Enumerate library compounds from building blocks Generate all possible chemical structures by combining building blocks according to the defined reaction schemes: delt-hit library enumerate --config\_path=/path/to/config.yaml --graph\_only=True  \# produce reaction graph only for inspection delt-hit library enumerate --config\_path=/path/to/config.yaml Those commands: Construct a reaction graph from your building blocks and reactions Enumerates all valid compound combinations Validates chemical structures Saves the complete library catalog  Expected outputs: library.parquet: Complete compound catalog with SMILES and barcode mappings reaction\_graph.png: Visual representation of synthetic schema additional\_reactions\_graph.png: Visualizes the additional reactions from the reaction\_graph sheet building\_block\_reactions\_graph.png: Visualizes the reactions derived from the building block sheets B0, B1, etc.  Figure 2 \textbar{} Complete reaction graph for a DEL with two building blocks (violet), reactions (green), products (red) and starting compounds (blue).   Performance: Enumeration speed is approximately 1’000-10’000 compounds/second depending on reaction complexity. For a 1 million compound library, expect 2-10 minutes of processing time depending on the machine. TROUBLESHOOTING: Check that all reaction SMIRKS are valid and specific enough to produce a single product Ensure reactant and product names match between building blocks and compounds sheets Re-run enumeration in debugging mode to visualize the reaction graphs that fail: delt-hit library enumerate --config\_path=/path/to/config.yaml --debug=invalid  Step 5 \textbar{} Calculate molecular properties and descriptors Compute comprehensive molecular descriptors for drug-likeness assessment and chemical space characterization: delt-hit library properties --config\_path=/path/to/config.yaml Computed properties include: Lipinski descriptors: Molecular weight (MW), partition coefficient (LogP), hydrogen bond donors (HBD), hydrogen bond acceptors (HBA) Topological indices: Topological polar surface area (TPSA), rotatable bonds Structural complexity metrics: Number of aromatic rings, heavy atom count, formal charge Drug-likeness scores: Quantitative estimate of drug-likeness (QED), fraction of sp³ carbons Expected outputs: properties/properties.parquet: Library catalog augmented with molecular descriptors properties/prop\_*.png: Distribution plots for each calculated property  Figure 3 \textbar{} Exemplary plot for the molecular weight distribution of library members.  Performance note: Property calculation processes approximately 600-1000 compounds/second. For a 1 million compound library, expect 15-30 minutes processing time. Step 6 \textbar{} Generate molecular representations for machine learning Create standardized chemical representations compatible with machine learning frameworks: \# Morgan fingerprints for similarity analysis \# \textasciitilde\{\}600-1000 compounds/s delt-hit library represent --method=morgan --config\_path=/path/to/config.yaml  \# BERT embeddings for deep learning applications \# \textasciitilde\{\}600-1000 compounds/s delt-hit library represent --method=bert --config\_path=/path/to/config.yaml Available representation methods: morgan: Morgan circular fingerprints (2048-bit vectors) for similarity searching and clustering bert: Transformer-based molecular embeddings (768-dimensional vectors) for deep learning Expected outputs: representations/morgan.npz: Sparse matrix of Morgan fingerprints representations/bert.npz: Dense matrix of BERT embeddings Note: These representations are compatible with scikit-learn, PyTorch, and TensorFlow for downstream machine learning applications. \\
\bottomrule
\end{tabular}
\end{table}